\section{The block graph transform}
\label{app:graph-transform}

We record the elementary similarity lemma with explicit constants.  Let
$K$ be a Hilbert space and
\[
  T=\begin{pmatrix}\lambda&\varphi\\0&0\end{pmatrix}
  \quad\text{on }\C\oplus K,
  \qquad \lambda\ne0,
\]
where $\varphi:K\to\C$.  Its spectral idempotent at $\lambda$ is
\[
  P=\lambda^{-1}T
   =\begin{pmatrix}1&g\\0&0\end{pmatrix},
  \qquad g=\lambda^{-1}\varphi.
\]
The row-operator norm gives
\begin{equation}
  \|P\|=(1+\|g\|^2)^{1/2}.
  \label{eq:projection-graph-norm}
\end{equation}
Define
\[
  Y=\begin{pmatrix}1&-g\\0&I\end{pmatrix},
  \qquad
  Y^{-1}=\begin{pmatrix}1&g\\0&I\end{pmatrix}.
\]
Then
\begin{equation}
  Y^{-1}TY=\begin{pmatrix}\lambda&0\\0&0\end{pmatrix},
  \qquad
  \|Y\|,\|Y^{-1}\|\le1+\|g\|.
  \label{eq:graph-transform-bound}
\end{equation}
Thus a uniform bound on the block idempotents is precisely what is needed
for the direct sum of the graph transforms to be bounded and invertible.

Conversely, if $T=XNX^{-1}$ and $N$ is compact normal, the Riesz
projections $Q_m$ of $N$ at distinct nonzero eigenvalues are orthogonal.
The corresponding projections of $T$ are $XQ_mX^{-1}$, so
\[
  \sup_m\|XQ_mX^{-1}\|\le\|X\|\|X^{-1}\|.
\]
This proves the necessity and sufficiency used in
\cref{lem:uniform-diagonalization} without appealing to a finite-dimensional
condition number that may depend on the block.

